\section{Correctness of the higher-order interpretation}
\label{sec:conservativity}

\Replaced{Two facts certify that the higher-order interpretation behaves correctly at first-order
types: every term of first-order type denotes, up to the isomorphisms relating the two
interpretations, a morphism of the first-order category; and the underlying function of the
interpretation agrees with the plain set-level interpretation of the term. The first is immediate,
because the embedding of the first-order model into the higher-order one is full
(\remref{conservativity:full-F}). The second is
proved by glueing: the interpreting category embeds into a category of Grothendieck logical
relations \citep{fiore-simpson99}, the language is interpreted there, and at first-order types the
interpretation is isomorphic to an embedded object, so by fullness of the embedding every term
denotes the image of a morphism of the embedded category (\thmref{conservativity:underlying}).}{The interpretation of the
language in the higher-order model is \emph{conservative} over the first-order model: every term
whose context and type are first-order denotes, up to the isomorphisms relating the two
interpretations, a morphism of the first-order category. The proof is by glueing: the first-order
category embeds into a category of Grothendieck logical relations \citep{fiore-simpson99} over the
higher-order one; the language is interpreted there; and at first-order types the interpretation
collapses back to the first-order category.}

Glueing along a functor is the standard method for definability results, also known as
sconing \citep{mitchell-scedrov92}; \citet{fiore-simpson99} refine it with \emph{covers},
decompositions of an object as a coproduct (\secref{conservativity:glueing}), which sum types make
necessary. \secref{predicate-system} sets up the calculus of predicates the construction uses,
\secref{conservativity:glueing} builds the glued category, and \secref{conservativity:definability}
proves the embedding full and derives the \Replaced{agreement of the underlying
interpretation}{conservativity theorem}. The two remaining subsections
supply the ingredients specific to inductive types: $\Poly$-types for the glued category
(\secref{conservativity:glr-poly-types}) and their preservation by the embedding
(\secref{conservativity:gf-poly-types}), which the theorem assumes.

\subsection{Predicate systems}
\label{sec:predicate-system}

\subsubsection{Predicate system over $\cat{C}$}

Let $\cat{C}$ be a category with finite products. A \emph{predicate system} over $\cat{C}$ consists of the
following data:
\begin{itemize}
\item For every object $X$ of $\cat{C}$, a poset $(\Pred(X), \sqsubseteq)$ of predicates on $X$.
\item For every $X$, the structure of a Heyting algebra (without $\bot$) on $\Pred(X)$, i.e.~ top element
   $\top$, meets $\meet$, joins $\join$ and residual $\residual$ satisfying the usual laws.
\item For every morphism $f: X \to Y$ of $\cat{C}$, a Galois connection $\directImage{-}{f} \dashv
\reindex{-}{f} : \Pred(Y) \to \Pred(X)$. The right adjoint $\reindex{-}{f}: \Pred(Y) \to \Pred(X)$ computes
the pullback (reindexing) of a predicate along $f$; the left adjoint $\directImage{-}{f}: \Pred(X) \to
\Pred(Y)$ computes the pushforward (existential image) of a predicate along $f$, with the adjointness property
meaning that
\[\directImage{\varphi}{f} \sqsubseteq \psi \iff \varphi \sqsubseteq \reindex{\psi}{f} \]
\item For every $X$ and $Y$, a right adjoint $\forall_Y: \Pred(X \times Y) \to \Pred(X)$ to the monotone map
$\reindex{-}{\pi_1}$ called \emph{universal quantification}.
\item $\reindex{-}{f}$ is colax (and thus $\directImage{-}{f}$ is lax) with respect to $\top$, $\meet$,
$\residual$ and $\forall_Y$.
\item $\reindex{-}{f}$ is lax (and thus $\directImage{-}{f}$ is colax) with respect to $\join$.
\end{itemize}

\noindent
A predicate system is a first-order hyperdoctrine, with posets of predicates in place of
categories: predicates are indexed by the objects of $\cat{C}$, substitution along a morphism acts
on them by reindexing, and the quantifiers are the adjoints to substitution, existential on the
left and universal on the right, which is the shape of the definition above. The reading of
quantifiers as adjoints originates with \citet{lawvere69}; \citet{pitts01} surveys the resulting
notion and its laws.

\subsubsection{Predicate system over $\Set$}

Let $X$ be a set. Define a \emph{predicate on $X$} to be a family of (proof-irrelevant) propositions for
every $x \in X$. For any function $f: X \to Y$, define $\reindex{\varphi}{f}(x)$ iff $\varphi(f(x))$ and
$\directImage{\varphi}{f}(y)$ iff there exists $x$ such that $f(x) = y$ and $\varphi(x)$. Define $\mathcal{S}$ to be the
predicate system over $\Set$ which assigns to any set $X$ the set $\Pred(X)$ of propositions on $X$ ordered by
implication, with $\top, \meet, \join$ and $\residual$ defined pointwise and universal quantification defined
as $\forall_Y(\varphi)(x)$ iff $\varphi(x,y)$ for all $y \in Y$.

\subsubsection{Predicates on a presheaf}

Let $\cat{C}$ be a category and $H: \cat{C}^\op \to \Set$ a presheaf. Define a \emph{predicate on $H$} to be
a family of predicates $\varphi_X \in \Pred_\mathcal{S}(H(X))$ for every object $X$ in $\cat{C}$ such that $\varphi_Y
\sqsubseteq \reindex{\varphi_X}{H(f)}$ for every $f: X \to Y$ in $\cat{C}$.

Now define the predicate system $\mathcal{S}^{\cat{C}}$ over $\Func{\cat{C}^\op}{\Set}$ which assigns to every
presheaf $H$ the poset $\Pred(H)$ of predicates on $H$ ordered pointwise, where $\top, \meet, \join,
\residual$ and $\forall$ in $\Pred(H)$ are all given pointwise, and to every natural transformation $\theta: H
\to K$, families of pullback and pushforward maps where $\reindex{\varphi}{\theta}_X = \reindex{\varphi_X}{\theta_X}$ and
$\directImage{\varphi}{\theta}_X = \directImage{\varphi_X}{\theta_X}$.

\subsubsection{Closure operator over a predicate system}

A \emph{closure operator} $C$ over a poset $(L, \sqleq)$ is a lax idempotent monad on $L$, with functoriality,
unit and multiplication given by monotonicity, extensivity ($p \sqleq C(p)$) and idempotence ($C(C(p)) \sqleq
C(p)$).

Let $\mathcal{S} = (\Pred, \top, \meet, \join, \residual, \forall)$ be a predicate system over a category
$\cat{C}$. A closure operator $C$ over $\mathcal{S}$ is a fibred lax idempotent monad $C$ on the fibration
$\Pred: \cat{C}^\op \to \Poset$ where for every object $X$ the fibre map $C_X: \Pred(X) \to \Pred(X)$ is a
closure operator, $C$ is lax monoidal with respect to $\meet$ (i.e.~$C(\varphi) \meet C(\psi) \sqleq C(\varphi \meet \psi)$),
and $C$ is lax strong, i.e.~$C(\varphi) \meet \psi \sqleq C(\varphi \meet \psi)$. $C$ being a fibred functor means it commutes
with any reindexing functor $\reindex{-}{f}$, i.e.~$C_X \comp \reindex{-}{f} \iso \reindex{-}{f} \comp C_Y$.

\subsection{The category of logical relations}
\label{sec:conservativity:glueing}

Fix a category $\cat{C}$ with terminal object, finite products and stable set-indexed coproducts,
in which the first-order fragment of the language is interpreted; a category $\cat{D}$ with terminal
object, finite products, exponentials, and set-indexed coproducts, in which the full
language is interpreted; and a functor $F : \cat{C} \to \cat{D}$ preserving the terminal object,
products, and set-indexed coproducts. The finite coproducts both interpretations need for sum types
are the two-element instance of the set-indexed ones, so they are not assumed separately; likewise
their stability and preservation by $F$ follow from the set-indexed versions. For the
language without inductive types it suffices to assume stable \emph{finite} coproducts, as in the
extensive categories of \citet{fiore-simpson99}; recursive types force the set-indexed
strengthening (\secref{conservativity:gf-poly-types}). Call a set-indexed coproduct decomposition
$\coprod_i y_i \cong y$ a \emph{cover} of the object $y$; covers generate the closure operator
introduced below.

\begin{lemma}
\label{lem:conservativity:fam-stable-coproducts}
In a $\Fam$ category the set-indexed coproducts are stable, and a functor acting fibrewise
preserves them.
\end{lemma}

\begin{proof}[Proof sketch]
A cover $\coprod_i y_i \cong y$ presents each index of $y$ as lying in a unique summand; pulling
back along $g$ retags the indices of the source, splitting it over the same index set, which gives
stability. A fibrewise functor leaves index sets untouched and acts on fibres alone, so it carries
the disjoint union of index sets to the disjoint union of their images.
\end{proof}

Predicates live over presheaves on $\cat{C}$: let $G : \cat{D} \to \PSh(\cat{C})$ be the functor
sending $X$ to the presheaf $\cat{D}(F(-), X)$. A presheaf predicate on $G(X)$ thus
assigns to each $y \in \cat{C}$, called the \emph{stage}, a predicate on the morphisms $F(y) \to X$, compatibly with
precomposition along $\cat{C}$-morphisms; the predicate system of \secref{predicate-system} provides
reindexing and the connectives for such predicates.

\begin{definition}[Definability]
\label{def:conservativity:definable}
For $x \in \cat{C}$, the predicate $\Definable_x$ on $G(F(x))$ holds at $y \in \cat{C}$ of those $f :
F(y) \to F(x)$ that are images of first-order morphisms: those with $f = F(g)$ for some $g : y \to x$.
\end{definition}

Definability is not closed under the reasoning the interpretation requires; in particular, a morphism
defined by case analysis on a coproduct is definable only up to precomposition with the case split.
Predicates are therefore taken \emph{closed} with respect to a closure operator $\mathbf{C}$,
generated by covers of the stage: $\mathbf{C}(\Phi)$ is defined
inductively, holding at $f : F(y) \to X$ if $\Phi$ does, or if some cover of $y$ has
$\mathbf{C}(\Phi)$ holding at each restriction of $f$.

\begin{lemma}
\label{lem:conservativity:indexed-closure}
$\mathbf{C}$ is a strong closure operator on predicates (\secref{predicate-system}), its
reindexing law using the stability of $\cat{C}$'s set-indexed coproducts.
\end{lemma}

\begin{proof}[Proof sketch]
Monotonicity, extensivity, idempotence, and strength are structural recursions on cover trees.
The reindexing law transports a cover tree along a morphism by pulling each cover back
(\lemref{conservativity:fam-stable-coproducts}) and recurring on the summands.
\end{proof}

The category $\GLR(F)$ of \emph{Grothendieck logical relations} has as objects pairs $(X, \Phi)$ of an
object $X \in \cat{D}$ and a closed predicate $\Phi$ on $G(X)$, and as morphisms $(X, \Phi) \to (Y,
\Psi)$ those $f : X \to Y$ with $\Phi \sqsubseteq \Psi[G(f)]$. \Replaced{The first-order category
embeds by}{Theorem~5.1 of the paper, after \citet{fiore-simpson99},
summarises the structure this category carries: $\GLR(F)$ is bicartesian closed
(here moreover with set-indexed coproducts), each piece built on the corresponding structure of
$\cat{D}$ with an appropriate predicate component; the projection $p : \GLR(F) \to \cat{D}$, $(X,
\Phi) \mapsto X$, strictly preserves it; and the first-order category embeds,}
\[
  \GF : \cat{C} \to \GLR(F) \qquad \GF(x) = (F(x), \mathbf{C}(\Definable_x)) \qquad \GF(g) = F(g)
\]
\Added{and the whole construction is summarised by the following theorem.}

\begin{AddedBlock}
\begin{theorem}[after \citet{fiore-simpson99}]
\label{thm:conservativity:glr}
$\GLR(F)$ is bicartesian closed, with set-indexed coproducts, each deriving from the
corresponding structure of $\cat{D}$ with an appropriate predicate component. The projection $p :
\GLR(F) \to \cat{D}$, $(X, \Phi) \mapsto X$, strictly preserves this structure, and $p \circ \GF =
F$. The embedding $\GF$ is full and preserves the terminal object, products and coproducts.
\end{theorem}

\begin{remark}
\label{rem:conservativity:glr-deltas}
Relative to the statement of \citet{fiore-simpson99}: $\cat{C}$ need not be cartesian closed, nor
$F$ preserve exponentials (when they are and it does, $\GF$ preserves exponentials too, but nothing
below needs this); the smallness of $\cat{C}$ is replaced, in the formalisation, by an appropriate
universe level; and the covers generating the closure operator are set-indexed rather than binary,
a strengthening required for inductive types (\secref{conservativity:gf-poly-types}). Enlarging the covers
strengthens the closure condition, so fewer predicates are closed and the category itself changes;
the theorem carries over because its proof uses only the closure laws of
\secref{predicate-system}, not the choice of generating covers.
\end{remark}
\end{AddedBlock}
\noindent
The functors assemble as
\begin{center}
\begin{tikzcd}[row sep=2.4em, column sep=3.2em]
  \cat{C} \arrow[r, "\GF"] \arrow[dr, "F"'] & \GLR(F) \arrow[d, "p"] & \\
  & \cat{D} \arrow[r, "G"'] & \PSh(\cat{C})
\end{tikzcd}
\end{center}
\noindent
with $p \circ \GF = F$: an object of $\GLR(F)$ is a $\cat{D}$-object together with a closed
predicate on the presheaf that $G$ assigns to it, and $\GF$ equips each $F(x)$ with the closure of
its definability predicate.

\begin{DeletedBlock}
For the language without inductive types, binary case splits, the two-element instance
of the covers above, suffice to generate the closure operator. Enlarging the class of covers
strengthens the closure condition, so fewer predicates are closed and the two glued categories
differ. The structure summarised above nevertheless carries over unchanged, because its proofs use
only the closure laws of \secref{predicate-system} and not the choice of generating
covers.
\end{DeletedBlock}

\subsection{Definability}
\label{sec:conservativity:definability}

The glueing construction makes $\GF$ \emph{full}: morphisms between embedded objects are
already first-order. Fullness is \Replaced{the last clause of \thmref{conservativity:glr}; here we
give the proof}{asserted as part of Theorem~5.1 of the paper; here we give the proof}.

\begin{lemma}[Fullness of $\GF$]
\label{lem:conservativity:definability}
Every morphism $f : \GF(x) \to \GF(y)$ of $\GLR(F)$ is the image of a first-order morphism: $f = F(g)$
for some $g : x \to y$.
\end{lemma}

\begin{proof}
As a morphism of $\GLR$, $f$ maps $\mathbf{C}(\Definable_x)$ into $\mathbf{C}(\Definable_y)$ reindexed along
$G(f)$; concretely, postcomposition with $f$ sends morphisms satisfying $\mathbf{C}(\Definable_x)$ to
morphisms satisfying $\mathbf{C}(\Definable_y)$. The identity $F(\id_x)$ satisfies $\Definable_x$, so $f = f
\circ F(\id_x)$ satisfies $\mathbf{C}(\Definable_y)$. But $\Definable_y$ is itself closed, because, node by
node in a cover tree, witnesses for the summand restrictions copair along the decomposition to a witness
for the covered morphism. So $f$ satisfies $\Definable_y$, that is, $f = F(g)$ for some $g : x \to y$.
\end{proof}

The copairing chooses a witness for each summand. Classically this is an instance of the axiom
of choice; in the formalisation, where definability is a proof-irrelevant existential, no choice
principle is available, and finite covers are special only because finitely many witnesses can
be extracted one at a time. We therefore assume a uniform choice of witnesses: a function that,
for each morphism $h : F(a) \to F(b)$ in the image of $F$, yields a $g$ with $F(g) = h$. When
$F$ is a fibrewise functor between $\Fam$ categories and is faithful, the assumption reduces to
its fibre functor: the index part of a witness is determined by $h$, the fibre parts are given
by the fibre functor's witness function, and faithfulness supplies the naturality of the
assembled family. \Replaced{At the pairing functor of \thmref{conservativity:underlying} below, a
morphism in the image is a pair whose second component is determined by the first, so the witness
is the first component itself.}{At the current instances the fibre functor is full by construction,
morphisms of the self-dual category being morphisms of the underlying one, so the witness is $h$
itself.}

Fullness of $\GF$ (\lemref{conservativity:definability}) applies only to morphisms between embedded
objects, so it remains to show that the interpretation of every first-order type is isomorphic to
one. Interpreting the language in $\GLR(F)$ needs, beyond the structure of
\secref{conservativity:glueing}, only $\Poly$-types (\defref{polynomial-types:has-poly}), which we
construct in \secref{conservativity:glr-poly-types}; their preservation by $\GF$ is assumed here and
proved in \secref{conservativity:gf-poly-types}. First-order types are also
interpreted directly in $\cat{C}$, which has $\Poly$-types whenever it is a $\Fam$ category
(\propref{polynomial-types:fam-has-poly}); write $\sem{\tau}_{\mathit{fo}}$ for this interpretation.

\begin{AddedBlock}
\begin{lemma}[Agreement]
\label{lem:conservativity:comparison}
For first-order $\tau$, $\sem{\tau} \cong \GF(\sem{\tau}_{\mathit{fo}})$.
\end{lemma}

\begin{proof}
By induction on $\tau$. The unit, product and sum cases follow from $\GF$'s preservation of the
terminal object, products and coproducts; base types agree by construction, the $\GLR(F)$-model of
the signature being the transport of the $\cat{C}$-model along $\GF$; and $\mu$-types follow from
the assumed preservation together with the isomorphism-invariance of the $\mu$-operator
(\secref{polynomial-types:fam}).
\end{proof}
\end{AddedBlock}

\begin{AddedBlock}
\begin{remark}[Full embeddings]
\label{rem:conservativity:full-F}
When $F$ is full (for example the embedding $\Fam(\SDSemiMod(S)) \to \Fam(\SemiMod(S))$), every
$\cat{D}$-morphism between embedded objects is an $F$-image outright, so conservativity over
$\cat{C}$ follows from the isomorphisms $\sem{\tau}_{\cat{D}} \cong F(\sem{\tau}_{\mathit{fo}})$
alone, with no glueing; these arise as in \lemref{conservativity:comparison}, with
$F$ in place of $\GF$.
\end{remark}
\end{AddedBlock}

\Added{The glueing machinery is needed for the following theorem, whose pairing functor
$\langle \mathrm{Id}, U \rangle$ is not full: the language is interpreted twice, in the
higher-order category and in a set-level category, and fullness of the glued embedding ensures
the two agree.}

\begin{AddedBlock}
\begin{theorem}[Underlying interpretation]
\label{thm:conservativity:underlying}
Suppose the set-indexed coproducts of $\cat{D}$ are stable. Let $\cat{E}$ be a further category with the
structure of $\cat{D}$ (\secref{conservativity:glueing}) and $\Poly$-types, and $U : \cat{D} \to
\cat{E}$ a functor preserving the terminal object, products and set-indexed coproducts.
Instantiate \secref{conservativity:glueing} with the first-order category taken to be $\cat{D}$
itself, the higher-order category $\cat{D} \times \cat{E}$, and
$F = \langle \mathrm{Id}, U \rangle$, and assume $\GF$ preserves $\Poly$-types. Then for every term
$\Gamma \vdash t : \tau$ with $\Gamma$ and $\tau$ first-order,
\[
  U(\sem{t}_{\cat{D}}) \;=\; \sem{t}_{\cat{E}},
\]
modulo the isomorphisms of \lemref{conservativity:comparison}.
\end{theorem}

\begin{proof}
A morphism $F(x) \to F(y)$ of $\cat{D} \times \cat{E}$ is a pair $(h, k)$ with $h$ and $k$
unrelated, whereas an $F$-image is constrained to $(g, U(g))$; so $F$ is not full and
\lemref{conservativity:definability} does real work. The interpretation of the language in
$\cat{D} \times \cat{E}$ is computed componentwise, so it is the pair $(\sem{t}_{\cat{D}},
\sem{t}_{\cat{E}})$. At first-order $\Gamma$ and $\tau$ the isomorphisms of
\lemref{conservativity:comparison} exhibit this
pair as a morphism between $\GF$-images, so by \lemref{conservativity:definability} it equals
$(g, U(g))$ for some $g$; the first component forces $g = \sem{t}_{\cat{D}}$, and the second then
reads $U(\sem{t}_{\cat{D}}) = \sem{t}_{\cat{E}}$.
\end{proof}
\end{AddedBlock}

\Added{For the canonical models, $U : \Fam(\SemiMod(S)) \to \Set$ sends a family to its index set
and a morphism to its index function; the stability hypothesis holds in any $\Fam$ category
(\lemref{conservativity:fam-stable-coproducts}). The theorem then says that the underlying function of the
higher-order interpretation is the plain set-theoretic interpretation of the term: the dependency
data does not interfere. $U$ does not preserve exponentials, which is why
\remref{conservativity:glr-deltas}'s weakening of the hypotheses on $F$ matters.}

\subsection{$\Poly$-types in $\GLR(F)$}
\label{sec:conservativity:glr-poly-types}

$\GLR(F)$ has $\Poly$-types. We prove this by transferring the construction of
\secref{polynomial-types:fam}: realise families of $\GLR(F)$-objects as set-indexed coproducts, and
obtain $\Poly$-types for $\GLR(F)$ from those of $\Fam(\GLR(F))$. The construction works for
any category with the structure listed in \thmref{conservativity:transfer}; preservation by $\GF$
is the subject of \secref{conservativity:gf-poly-types}.

\paragraph{Realisation.}
Let $\cat{E}$ be a category with set-indexed coproducts. The functor
\[
  \Sigma : \Fam(\cat{E}) \to \cat{E} \qquad \Sigma(X, \partial X) = \coprod_{x \in X} \partial X_x
\]
sends a family to the coproduct of its fibres and a morphism to the copairing of its fibre maps (the
counit of the free coproduct completion). A \emph{presentation} of an object $Z \in \cat{E}$ is a
family together with an isomorphism $\Sigma(X, \partial X) \cong Z$. In a $\Fam$ category every
object is canonically presented, set-indexed coproducts being disjoint unions of index sets; in
$\GLR(F)$ presentations are genuine data, since a closed predicate on the carrier of $Z$ need not
be determined by its restrictions to the fibres of a presenting family.

\paragraph{$\Poly$-types via realisation.}
A polynomial $P$ over $\cat{E}$ embeds as a polynomial $\hat{P}$ over $\Fam(\cat{E})$, each
$\const$-object replaced by its $\eta$-image, and $\Fam(\cat{E})$ has $\Poly$-types by
\propref{polynomial-types:fam-has-poly}, since $\cat{E}$ has a terminal object and finite
products. The $\Poly$-types for $\cat{E}$ are the realisations: $\mu_{\alpha.P}(\delta) :=
\Sigma\, \mu_{\alpha.\hat{P}}(\eta\,\delta)$, that is, the composite
\begin{center}
\begin{tikzcd}[row sep=2.4em, column sep=4.5em]
  \Fam(\cat{E})^{\Delta} \arrow[r, "\mu_{\alpha.\hat{P}}"] & \Fam(\cat{E}) \arrow[d, "\Sigma"] \\
  \cat{E}^{\Delta} \arrow[u, "\eta^{\Delta}"] \arrow[r, dashed, "\mu_{\alpha.P}"'] & \cat{E}
\end{tikzcd}
\end{center}
\noindent
which routes an environment through $\Fam(\cat{E})$, where the initial algebras exist by
construction, and realises the result back in $\cat{E}$. For this to yield
\defref{polynomial-types:has-poly}, $\Sigma$ must preserve what the construction uses:
set-indexed coproducts (immediate), the terminal object, and finite products, the last requiring
products to distribute over set-indexed coproducts (a consequence of the exponentials).

\paragraph{Operations and laws.}
These do not simply transport along the resulting isomorphisms: the algebra map of
$\mu_{\alpha.\hat{P}}$ lives at an environment containing the $\Fam(\cat{E})$-object
$\mu_{\alpha.\hat{P}}$ itself, whereas its realisation must live at the presented object
$\Sigma\, \mu_{\alpha.\hat{P}}$. The adjunction $\Sigma \dashv \eta$, where $\eta : \cat{E}
\to \Fam(\cat{E})$ embeds an object as a family over a one-element set, transposes algebras and
catamorphisms between the two categories; what remains is that realisation is invariant under
replacing environment entries by families with isomorphic realisations. We package this as a
property of the polynomial, established by induction.

\begin{definition}[Collapse]
\label{def:conservativity:collapse}
A polynomial $P$ over $\cat{E}$ \emph{collapses} if every family of isomorphisms $\Sigma\,
\hat{\delta}_i \cong \Sigma\, \hat{\delta}'_i$ between the realisations of two environments
induces an isomorphism $\Sigma(\hat{P}(\hat{\delta})) \cong \Sigma(\hat{P}(\hat{\delta}'))$,
compatibly with composition of such families, and naturally with respect to the strong action of
$\hat{P}$.
\end{definition}

\begin{remark}
Compatibility with identities is derivable: naturality at projection families gives an
extensionality principle (collapse depends only on the underlying isomorphisms), extensionality
makes the collapse at identity families idempotent via compatibility with composition, and an
idempotent isomorphism is the identity. Compatibility with composition is not derivable from
naturality: the naturality squares require morphisms of $\Fam(\cat{E})$ along the environment
directions, and none exist along a family of mere isomorphisms of realisations.
\end{remark}

\begin{lemma}
\label{lem:conservativity:collapse}
Every polynomial collapses. For $\mu$-polynomials the induced isomorphism is a morphism of
algebras, and the laws of \defref{conservativity:collapse} follow from uniqueness of
catamorphisms.
\end{lemma}

\begin{theorem}[Transfer]
\label{thm:conservativity:transfer}
Let $\cat{E}$ have set-indexed coproducts, a terminal object, finite products, exponentials and
strong coproducts. Then $\cat{E}$ has $\Poly$-types (\defref{polynomial-types:has-poly}).
\end{theorem}

\begin{proof}[Proof sketch]
Take $\mu_{\alpha.P} := \Sigma\, \mu_{\alpha.\hat{P}}$ at singleton environments. The algebra
map is the realised $\Fam(\cat{E})$ algebra map corrected by collapse
(\lemref{conservativity:collapse}) at the bound entry; the catamorphism transposes the
$\Fam(\cat{E})$ catamorphism along $\Sigma \dashv \eta$. For the $\beta$ and $\eta$ laws, the
interpretations agree across a comparison isomorphism $\Sigma(\hat{P}(\hat{\delta})) \cong
P(\delta)$, and the realised strong action of $\hat{P}$ simulates the strong action that $\cat{E}$
derives from the new structure, naturally in this comparison (induction on $P$; the $\mu$ case is
naturality of collapse). The laws are then the $\Fam(\cat{E})$ laws conjugated by the simulation.
\end{proof}

\begin{corollary}
\label{cor:conservativity:glr-poly}
$\GLR(F)$ has $\Poly$-types, its structure being that of
\secref{conservativity:glueing}, with set-indexed coproducts computed as coproducts of carriers
with joined predicates.
\end{corollary}

\subsection{Preservation of $\Poly$-types by $\GF$}
\label{sec:conservativity:gf-poly-types}

Preservation does not follow from $\GF$'s preservation of the finite structure: initial algebras
map out of themselves, so finite preservation yields only the comparison $\mu_{\alpha.\GF(P)}
\to \GF(\mu_{\alpha.P})$, by folding the $\GF$-image of the algebra. The converse direction
requires $\GF$ to preserve the presentation of $\mu_{\alpha.P}$ as a set-indexed coproduct of its
fibres.

Preservation asks only for an isomorphism of objects, so the comparison can be carried out on the
concrete carriers of \propref{polynomial-types:fam-has-poly}, taking $\cat{C}$ to be a $\Fam$
category with those $\Poly$-types; which $\Poly$-types structure is immaterial, any two having
componentwise isomorphic $\mu$-objects, by folds both ways and uniqueness of catamorphisms. On the
carriers no initiality machinery is needed: the tree sorts and fibres of the carrier are computed from the index sets and fibres of the
environment entries and $\const$-objects alone. In particular a $\const$-object and a variable
bound to an equal environment entry contribute identical data. For a polynomial $Q$, write $Q^{*}$
for its constant-free \emph{skeleton}: each $\const$-object replaced by a fresh variable, the
$\const$-objects $\vec{A}$ of $Q$ collected as additional environment entries. Having no constants,
$Q^{*}$ is a polynomial over any category; below it is interpreted over both $\Fam(\cat{C})$ and
$\Fam(\GLR(F))$, with the same sorts and trees on each side. Write
$\check{\delta}_i$ for the $\GF$-image of the canonical presentation of $\delta_i$: the family
over the index set of $\delta_i$ obtained by applying $\GF$ to the components of the
presentation.

\begin{lemma}
\label{lem:conservativity:definable-coproducts}
$\Definable_{\coprod_i x_i} \sqsubseteq \mathbf{C}(\bigvee_i \Definable_{x_i}\langle
\mathsf{in}_i \rangle)$.
\end{lemma}

\begin{proof}[Proof sketch]
A definable $f = F(g)$ into the coproduct decomposes along the pullback of the cover
$\coprod_i x_i$ along $g$, by stability; each restriction factors definably through a summand.
\end{proof}

Built-in list types would need the same extension. Together with preservation of set-indexed
coproducts by $F$ it makes $\GF$ preserve them: carriers by assumption, predicates by the lemma and
the closure operator's reindexing law (\lemref{conservativity:indexed-closure}), the coproducts of
$\GLR(F)$ being computed as in \corref{conservativity:glr-poly}. In particular
$\Sigma\check{\delta}_i \cong \GF(\delta_i)$.

\begin{lemma}[Skeleton]
\label{lem:conservativity:skeleton}
In a $\Fam$ category, $\mu_{\alpha.Q}(\delta) \cong \mu_{\alpha.Q^{*}}(\delta, \vec{A})$,
where $\vec{A}$ are the $\const$-objects of $Q$.
\end{lemma}

\begin{proof}[Proof sketch]
The carrier construction consumes a $\const$-object and the environment entry of a variable
through the same data: the index set, decorating the tree sorts, and the fibres, interpreting the
leaves. The constant/variable distinction is thus invisible to it, and the two carriers agree up to
the renaming of tree constructors induced by the skeleton, by tree recursion.
\end{proof}

\begin{lemma}[Carrier comparison]
\label{lem:conservativity:carrier-comparison}
$\GF(\mu_{\alpha.P^{*}}(\delta, \vec{A})) \cong \Sigma(\mu_{\alpha.P^{*}}(\check{\delta},
\check{A}))$.
\end{lemma}

\begin{proof}[Proof sketch]
Both carriers are set-indexed coproducts over the shared sort set, the left by the canonical
presentation of the $\Fam$ $\mu$-object and the right by realisation. Since $\GF$ preserves
set-indexed coproducts, the comparison reduces to the fibres over a common sort, by tree
recursion: at the leaves the two fibres are the same $\GF$-image of a presentation component, by
the construction of $\check{\delta}$ and $\check{A}$, and at products they agree by $\GF$'s
preservation of products.
\end{proof}

\begin{proposition}[Preservation]
\label{prop:conservativity:gf-preserves-poly}
$\GF$ preserves $\Poly$-types.
\end{proposition}

\begin{proof}[Proof sketch]
Apply \lemref{conservativity:skeleton} in $\cat{C}$, then
\lemref{conservativity:carrier-comparison}, then collapse (\lemref{conservativity:collapse}) at
the constant-free $P^{*}$, at the isomorphism family given by $\Sigma\check{\delta}_i \cong
\GF(\delta_i) \cong \Sigma(\eta\,\GF\,\delta_i)$ and likewise at the entries for
$\vec{A}$, and finally the realisation of \lemref{conservativity:skeleton} in $\Fam(\GLR(F))$,
backwards, at $\hat{P} = \eta\,\GF\,P$. The skeleton discards the $\const$-objects, so
$\hat{P}^{*} = P^{*}$, with the $\const$-objects of $\hat{P}$ the $\eta\,\GF$-images of
$\vec{A}$:
\[
  \GF(\mu_{\alpha.P})
  \;\cong\; \GF(\mu_{\alpha.P^{*}}(\delta, \vec{A}))
  \;\cong\; \Sigma(\mu_{\alpha.P^{*}}(\check{\delta}, \check{A}))
  \;\cong\; \Sigma(\mu_{\alpha.P^{*}}(\eta\,\GF\,\delta, \eta\,\GF\,\vec{A}))
  \;\cong\; \Sigma(\mu_{\alpha.\hat{P}}(\eta\,\GF\,\delta))
  \;=\; \mu_{\alpha.\GF(P)}.
\]
\end{proof}
